\documentclass[10pt,a4paper]{article}
\usepackage[utf8]{inputenc}
\usepackage[a4paper,left=3cm,right=2cm]{geometry}
\usepackage{amsmath}
\usepackage[thmmarks, amsmath, thref]{ntheorem}
\usepackage{amsfonts}
\usepackage{amssymb}
\usepackage{fancyhdr}
\pagestyle{fancy}
\usepackage{anysize}
\usepackage{graphicx}
\usepackage{float}
\usepackage{hyperref}
\usepackage{pgfplots}
\pgfplotsset{compat=1.17}
\usepackage{enumitem}
\newtheorem{theorem}{Theorem}
\theoremsymbol{\ensuremath{\blacksquare}}
\newtheorem*{solution}{Solución:}

\usepackage{xcolor}

% Margenes y formalidades

\textwidth 6.25in % Margen Derecho y ancho del texto
\textheight 8.5in
\oddsidemargin 0in % Margen Izquierdo
\headheight 0in % Largo del texto
\leftmargin 1in
\parindent 0em % Salto entre la indentación (Sangría)https://www.overleaf.com/4519127669sdpptkhvsbfn
\topmargin -0.5in % Margen Superior
\parskip 0.5ex % Salto entre parrafos
\baselineskip 1pt



\lhead{
\includegraphics[scale=0.11]{figuras/logoUSM.png}} % texto izquierda de la cabecera
%\chead{}                                      % texto centro de la cabecera
\rhead{Universidad Técnica Federico Santa María\\Ingeniería Comercial } % número de página a la derecha
\renewcommand{\headrulewidth}{0.4pt} % grosor de la línea de la cabecera
\renewcommand{\footrulewidth}{0.4pt} % grosor de la línea del pie

\begin{document}
\vspace{5cm}

\begin{center}
\begin{Large}
\textbf{\\Guía 1: Econometría ICS-294}
\end{Large}\\
\end{center}

\section{Muestreo}

\textbf{Ejercicio 1:} Se tiene una población de 1000 individuos con una media de ingreso mensual de $1500$ y una desviación estándar de $300$. Si se toma una muestra aleatoria simple de 100 individuos, ¿cuál es la media esperada y la desviación estándar de la media muestral?



\section{Estadística Descriptiva}

\textbf{Ejercicio 2:} Se tiene el siguiente conjunto de datos sobre el número de ventas realizadas por 5 vendedores en un mes:
$$X = \{10,15,12,18,20\}$$
Calcule la media, mediana y varianza muestral.


\section{Intervalos de confianza}

\textbf{Ejercicio 3:} Una empresa desea comparar los tiempos promedio de producción de dos plantas, Planta A y Planta B. Se toman muestras aleatorias de los tiempos de producción (en horas) de ambas plantas.

\begin{itemize}
    \item Para la \textbf{Planta A}, se selecciona una muestra de 50 observaciones con una media muestral de \( \bar{X}_A = 5,2 \) horas y una varianza conocida de \( \sigma_A^2 = 0,16 \) horas.
    \item Para la \textbf{Planta B}, se selecciona una muestra de 60 observaciones con una media muestral de \( \bar{X}_B = 4,9 \) horas y una varianza conocida de \( \sigma_B^2 = 0,25 \) horas.
\end{itemize}

 Verifique si la diferencia entre los tiempos de producción de las dos plantas es estadísticamente significativa.


\section{Regresión Lineal Simple}

\textbf{Ejercicio 4:} Se tiene la siguiente muestra de datos:
\begin{center}
\begin{tabular}{|c|c|}
\hline
X (Publicidad) & Y (Ventas) \\
\hline
2 & 4 \\
\hline
4 & 7 \\
\hline
6 & 10 \\
\hline
8 & 13 \\
\hline
10 & 16 \\
\hline
\end{tabular}
\end{center}

Ajuste un modelo de regresión lineal simple de la forma:
$$Y = \beta_0 + \beta_1 X + \epsilon $$

usando el método de mínimos cuadrados ordinarios.
\vspace{0.5cm}


\textbf{Ejercicio 5:} 
Se tiene un conjunto de datos que relaciona el número de horas de estudio (\(x\)) y la nota obtenida en un examen (\(y\)) por un grupo de 5 estudiantes:

\[
\begin{array}{|c|c|}
\hline
\text{Horas de estudio} (x) & \text{Nota} (y) \\
\hline
1 & 2 \\
2 & 3 \\
3 & 5 \\
4 & 4 \\
5 & 6 \\
\hline
\end{array}
\]

Se considera que la relación entre la nota obtenida y las horas de estudio 
est\'a descrita por el siguiente modelo regresión lineal poblacional:

\[
y = \beta_0 + \beta_1 x+u
\]

\begin{enumerate}
    \item  Usando el método de mínimos cuadrados, se pide calcular los estimadores de los parámetros \( \beta_0 \) y \( \beta_1 \). 

     \item Determine los valores predichos por el modelo para las horas de estudio de la muestra y sus  residuos. Verifique que los residuos suman 0.

     
    \item  Compruebe que la regresión pasa por las medias de las variables x e y (es decir, el punto $(\bar{x}, \bar{y})$ está en la recta).

\item Grafique la regresión lineal muestral, las observaciones de la muestra e identifique en gráfico los residuos.

    \item  Determine el $R^2$ de la regresión e interprete su resultado. 
\end{enumerate}

{\bf Ejercicio 6:} Sean \(\hat{\beta}_0\) y \(\hat{\beta}_1\) el intercepto y la pendiente en la regresión de \(y_i\) sobre \(x_i\), usando \(n\) observaciones.

Ahora, sean \(\tilde{\beta}_0\) y \(\tilde{\beta}_1\) la intersección con el eje \(y\) y la pendiente en la regresión de \((c_1 + y_i)\) sobre \((c_2 + x_i)\) (sin ninguna restricción sobre \(c_1\) o \(c_2\)). Muestre que \(\tilde{\beta}_1 = \hat{\beta}_1\) y \(\tilde{\beta}_0 = \hat{\beta}_0 +c_1- c_2 \cdot \hat{\beta}_1\). Compare los $R^2$ de ambas regresiones.
    
\section{Formulario}
\begin{itemize}

\item El intervalo de confianza de $100(1-\alpha)\%$ para la media poblacional $\mu$ cuando $\sigma$ es \emph{\textbf{conocido}}:

$$\color{red}{\bar{x} \pm z_{\frac{\alpha}{2}}\frac{\sigma}{\sqrt{n}}}$$

\item El intervalo de confianza de $100(1-\alpha)\%$ para la diferencia de medias
$\mu_x-\mu_y$ es:

$${\color{red}(\bar{X}-\bar{Y}) \pm z_{\frac{\alpha}{2}}\sqrt{\frac{\sigma_x^2}{n_x}+\frac{\sigma_y^2}{n_y}}}$$


\item El intervalo de confianza de $100(1-\alpha)\%$ para la media poblacional $\mu$ cuando $\sigma$ es \emph{\textbf{desconocido}}:

$$\color{red}{\bar{x} \pm t_{(n-1)\frac{\alpha}{2}}\frac{S}{\sqrt{n}}}$$

\item  Estimadores MCO.
{\color{red}
$$\hat{\beta}_0 = \bar{y} - \hat{\beta}_1 \bar{x} $$
$$\hat{\beta}_1 = \frac{\sum_{i=1}^{n}(x_i - \bar{x})(y_i-\bar{y})}{\sum_{i=1}^{n}(x_i - \bar{x})^2}$$}

    
\item SST la suma de cuadrados totales (total sum of squares):
	{\color{red}	\[ SST = \sum_{i=1}^{n} (y_i - \bar{y})^2 \]}
		
		\item  SSE la suma de cuadrados explicados (explained sum of squares):
	{\color{red}	\[ SSE = \sum_{i=1}^{n} (\hat{y}_i - \bar{y})^2 \]}
		
		\item  SSR la suma de cuadrados no explicados (residual sum of squares):
	{\color{red}	\[ SSR = \sum_{i=1}^{n} \hat{u}_i^2 \]}

\end{itemize}

\begin{figure}
    \centering
    \includegraphics[width=0.7\linewidth]{Figuras/normal_critic.png}
\end{figure}







\end{document}